\documentclass[../../main.tex]{subfiles}
\begin{document}

\paragraph{【新課程】}
点$P(x,y)$は$xy$平面上の円$C:(x-5)^2+(y-5)^2=r^2$ ($r>0$)の上を動く動点である．
このとき点$P$の点$A(9,0)$に関する対称点を$Q$とし，
また点$P$を原点$O$のまわりに正の向きに$\dfrac{\pi}{2}$だけ回転した点を$R$とする．
点$P$が円$C$の上を動くときの線分$QR$の長さの最小値$f(r)$と最大値$g(r)$とを求めよ．
また$f(r)$が$0$となるような$r$の値を求めよ．

\paragraph{【旧課程】}
$xy$平面上に3つの円$A$，$B$，$C$があって，それぞれ
\begin{align}
  A:x^2+y^2=9, \quad B:(x-4)^2+(y-3)^2=4, \quad C:(x-5)^2+(y+3)^2=1
\end{align}
で表わされる．
この平面上の点$P$から円$A$，$B$，$C$に接線がひけるとき，
$P$からそれらの接点までの距離をそれぞれ$\alpha(P)$，$\beta(P)$，$\gamma(P)$とする．
このとき${\alpha(P)}^2+{\beta(P)}^2+{\gamma(P)}^2=99$となる点$P$の全体が作る曲線を図示し，
その長さを求めよ．

\end{document}
